% from ega4-1-fr.tex
\oldpage[IV-1]{229}

\begin{proposition}[1.3.3]
\label{IV.1.3.3-fr}
Si $Y$ est un préschéma localement noethérien et $f:X\to Y$ un morphisme
localement de type fini, alors $X$ est localement noethérien.
\end{proposition}

Pour la démonstration, voir (I, 6.3.7).

\begin{proposition}[1.3.4]
\label{IV.1.3.4-fr}
\begin{enumerate}
  \item[(i)] Toute immersion locale est localement de type fini.
  \item[(ii)] Si deux morphismes $f:X\to Y$, $g:Y\to Z$ sont localement de
  type fini, il en est de même de $g\circ f$.
  \item[(iii)] Si $f:X\to Y$ est un $S$-morphisme localement de type fini,
  $f_{(S')}:X_{(S')}\to Y_{(S')}$ est localement de type fini pour toute
  extension $S'\to S$ du préschéma de base.
  \item[(iv)] Si $f:X\to X'$ et $g:Y\to Y'$ sont deux $S$-morphismes
  localement de type fini, $f\times_Sg$ est localement de type fini.
  \item[(v)] Si le composé $g\circ f$ de deux morphismes est localement de
  type fini, $f$ est localement de type fini.
  \item[(vi)] Si un morphisme $f$ est localement de type fini, il en est de
  même de $f_{\mathrm{red}}$.
\end{enumerate}
\end{proposition}

Pour la démonstration, voir (I, 6.6.6).

\begin{corollary}[1.3.5]
\label{IV.1.3.5-fr}
Soit $f:X\to Y$ un morphisme localement de type fini. Pour tout morphisme
$Y'\to Y$ tel que $Y'$ soit localement noethérien, $X\times_YY'$ est
localement noethérien.
\end{corollary}

\begin{proof}
En effet, $f_{(Y')}:X\times_YY'\to Y'$ est localement de type fini par
(1.3.4, (iii)), et il suffit d'appliquer (1.3.3).
\end{proof}

\begin{proposition}[1.3.6]
\label{IV.1.3.6-fr}
Soit $\varphi:A\to B$ un homomorphisme d'anneaux. Pour que le morphisme
correspondant $f:\operatorname{Spec}(B)\to\operatorname{Spec}(A)$ soit localement de type fini, il faut et
il suffit que $B$ soit une $A$-algèbre de type fini.
\end{proposition}

Pour la démonstration, voir (I, 6.3.3), en tenant compte de ce que $f$ est
quasi-compact et de (I, 6.6.3).

\begin{proposition}[1.3.7]
\label{IV.1.3.7-fr}
Pour qu'un morphisme localement de type fini $f:X\to Y$ soit surjectif, il
faut et il suffit que, pour tout corps algébriquement clos $\Omega$,
l'application $X(\Omega)\to Y(\Omega)$ correspondant à $f$ (I, 3.4.1) soit
surjective.
\end{proposition}

Pour la démonstration, voir (I, 6.3.10).

\begin{env}[1.3.8]
\label{IV.1.3.8-fr}
Soit $A$ un anneau. On dit qu'une $A$-algèbre $B$ est \emph{essentiellement de
type fini} si $B$ est $A$-isomorphe à une $A$-algèbre de la forme $S^{-1}C$,
où $C$ est une $A$-algèbre \emph{de type fini} et $S$ une partie multiplicative
de $C$.
\end{env}

\par\medskip\noindent
\begin{proposition}[1.3.9]
\label{IV.1.3.9-fr}
\begin{enumerate}
  \item[(i)] Si $B$ est une $A$-algèbre essentiellement de type fini et $C$
  une $B$-algèbre essentiellement de type fini, alors $C$ est une
  $A$-algèbre essentiellement de type fini.
  \item[(ii)] Soient $B$ une $A$-algèbre essentiellement de type fini, $A'$
  une $A$-algèbre~; alors $B'=B\otimes_AA'$ est une $A'$-algèbre
  essentiellement de type fini.
\end{enumerate}
\end{proposition}

\begin{proof}
\begin{enumerate}
  \item[(i)] Soient $B=S'^{-1}B'$ et $C=T'^{-1}C'$, où $B'$ (resp. $C'$) est
  une $A$-algèbre (resp. une $B$-algèbre) de type fini et $S'$ (resp. $T'$)
  une partie multiplicative de $B'$ (resp. $C'$)~; alors $C'$ est de la forme
  $S'^{-1}C''$, où $C''$ est une $B'$-algèbre de type fini, donc $C$ est aussi
  un anneau de fractions de $C''$~; comme $C''$ est une $A$-algèbre de type
  fini, cela prouve notre assertion.
  \item[(ii)] Si $B$ est anneau de fractions d'une $A$-algèbre de type fini
  $C$, $B'$ est un anneau de fractions de la $A'$-algèbre de type fini
  $C\otimes_AA'$, d'où (ii).
\end{enumerate}
\end{proof}

